% ============================================================
\chapter{Discussion and Research Outlook}
\label{ch:discussion}
% ============================================================

\section{Synthesis of Findings}

The PIDNet-S experiments establish that the CaT capability labels can be
learned by a compact network at useful accuracy and latency. The retained
7.62-million-parameter model reaches 0.893941 mIoU, and an independent CPU run
reconstructs its complete confusion matrix exactly. This level of
reproducibility connects the reported rates to one identifiable model rather
than to an undocumented training run. It remains a CaT result, however;
neither the exact score nor its operational interpretation transfers automatically
to another site or vehicle.

Several seemingly reasonable interventions did not work as intended. The
tested photometric augmentation reduced performance, the additional
false-safe term increased both asymmetric error rates, and greater input
resolution provided no benefit while halving the batch size. Neutral class
weights recovered more traversable pixels but did not give the strongest
combined operating point. By contrast, threshold adjustment behaved
predictably: it traded false-safe openings for false blocks without eliminating
the underlying errors. These outcomes argue for controlled measurements of
engineering choices rather than assumptions based on their names or intent.

Pretrained feature transfer changes the attainable accuracy, but training
horizon changes the comparison as well. Under the original matched 15-epoch
conditions, ROD improves mean mIoU over PIDNet-S by 3.15 percentage points for
blue+green and 3.79 points for blue-only. The convergence-aware blue+green
follow-up raises both means to 0.9331 and 0.9195, narrowing their difference to
1.36 points. Across all nine systems, FPN/EfficientNet-B0 reaches
$0.9423\pm0.0005$ and is highest in each of the three executed seeds;
U-Net/EfficientNet-B0 follows at $0.9409\pm0.0004$. ROD remains below the five
pretrained encoder--decoder alternatives but nearly matches
U-Net/MobileNetV2. These changes show that the earlier matrix described
performance after equal early exposure, not a converged architecture ranking.

Runtime supplies a different ordering. PIDNet-S produces the highest H100
throughput, compiled DDRNet-23-Slim leads on the tested CPU, and the two
EfficientNet-B0 systems provide the greatest overlap. TensorRT brings the
reconstructed ROD model to 39.38 FPS for the measured local pipeline, but
PIDNet-S remains much faster at 106.43 FPS. On the Ryzen CPU, ROD reaches only
0.41 eager FPS and 0.46 compiled FPS because freezing removes encoder gradients
during training but not encoder computation during inference. Model selection
is consequently a constraint problem involving target frame rate, runtime, and
acceptable error profile, not a lookup of the largest mIoU.

\section{Interpreting the Architecture Results}

The comparison does not support a simple narrative that transformers are more
accurate and CNNs are faster. ROD's pretrained transformer improves over
randomly initialised PIDNet-S, but FPN and U-Net with a compact pretrained
convolutional EfficientNet-B0 encoder perform better under the common recipe.
SegFormer-B0 lies between the EfficientNet and MobileNet systems. The best
observed result is therefore associated with strong pretrained compact
features and an effective multiscale decoder, not with one architecture
category alone.

The $2\times2$ subset suggests that encoder choice has the larger effect:
EfficientNet-B0 improves both decoders, while FPN's advantage over U-Net is
0.13 points with EfficientNet and 0.41 points with MobileNetV2. FPN is higher
for every matched encoder--seed pair in the convergence suite. With only two
encoders and two decoders, this remains a local interaction, not a universal
conclusion about FPN or U-Net.

ROD's reconstructed training result shows that optimisation is
architecture-dependent. CE alone, neutral weights, a larger learning rate, and
polynomial decay improve ROD relative to the common CE+Dice recipe. This does
not invalidate the retained PIDNet configuration. It shows that a recipe
selected during one architecture's development should not be assumed optimal
for another.

The real-time architectures occupy different parts of the speed--accuracy
plane. DDRNet provides the highest compiled CPU rate but the lowest
convergence-aware blue+green accuracy. PIDNet provides the highest H100
throughput, a verified deployment chain, and now a higher mean mIoU than
BiSeNetV2. The phrase ``real-time architecture'' therefore describes design
intent, not a fixed ranking on every processor.

\section{Meaning of the Asymmetric Error Metrics}

The reference PIDNet model has lower FSR than FBR, which indicates a conservative
aggregate operating balance. The maximum-FSR image demonstrates why that
statement is insufficient: a 4.32\% global rate can coexist with a 30.20\%
per-image rate and one connected expansion into a large field. Conversely, the
maximum-FBR example shows a visually coherent route being almost removed.

The highest-scoring models reduce both directions, so their accuracy gain is not
obtained simply by opening or closing the entire mask. Even then, pixel rates
do not specify whether an error intersects the planned trajectory. A false-safe
pixel behind the vehicle and one under the projected front wheel contribute
equally. Connectedness, distance, obstacle type, speed, and confidence are
absent.

The metrics are best used as three layers of evidence:
\begin{enumerate}
  \item mIoU and class IoU summarise balanced segmentation overlap;
  \item FSR and FBR expose the direction of disagreement;
  \item ranked visual review shows the spatial form of the error.
\end{enumerate}

Figure~\ref{fig:safety_evidence_ladder} shows why these layers should not be
collapsed into a single safety claim. This thesis evaluates pixel disagreement
and its spatial form. Whether an error changes a route, and whether the vehicle
then collides or loses progress, require a planner and physical system that are
outside the present experiment.

\begin{figure}[H]
\centering
\resizebox{\textwidth}{!}{\begin{tikzpicture}[
  font=\sffamily\small,
  stage/.style={draw,rounded corners=2pt,very thick,align=center,
    text width=2.75cm,minimum height=1.95cm,inner sep=5pt},
  future/.style={stage,dashed,draw=black!55,fill=black!5},
  flow/.style={-{Latex[length=2.4mm]},very thick,draw=azulUC3M}
]
  \node[stage,fill=softblue] (metrics) at (1.55,2.05)
    {\textbf{Pixel metrics}\\mIoU, FSR, FBR\\\footnotesize how many pixels disagree?};
  \node[stage,fill=softorange] (spatial) at (5.05,2.05)
    {\textbf{Spatial review}\\location, shape, connectedness\\\footnotesize what does the error look like?};
  \node[future] (planner) at (8.55,2.05)
    {\textbf{Planner interaction}\\footprint, route, distance\\\footnotesize does the error affect a path?};
  \node[future] (vehicle) at (12.05,2.05)
    {\textbf{Vehicle outcome}\\collision, intervention, progress\\\footnotesize what happened physically?};
  \draw[flow] (metrics) -- (spatial);
  \draw[flow] (spatial) -- (planner);
  \draw[flow] (planner) -- (vehicle);

  \draw[very thick,good] (0.05,-0.38) -- node[below,font=\footnotesize,text=black]
    {evaluated in this thesis} (6.55,-0.38);
  \draw[very thick,dashed,black!55] (7.05,-0.38) -- node[below,font=\footnotesize,text=black]
    {requires future system-level validation} (13.55,-0.38);
\end{tikzpicture}
}
\caption[From pixel metrics to vehicle-level evidence.]{Evidence required to
move from segmentation quality to a vehicle-level safety statement. Solid
stages are evaluated here; dashed stages require future planner and vehicle
experiments. Original schematic drawn for this thesis.}
\label{fig:safety_evidence_ladder}
\end{figure}

A later planner-level study should add collision, intervention, route
availability, progress, and time-to-goal measures rather than attempting to
turn FSR alone into a safety certificate.

\section{Annotation and Task Uncertainty}

CaT's vehicle-capability concept is a strength because it is closer to
navigation than a material label. It also introduces judgement. A sharp
semantic boundary between a track and adjacent vegetation may be visible in
RGB; the exact boundary between pickup-traversable and off-road-only soil may
depend on slope, compaction, moisture, hidden roughness, or an annotator's
vehicle model.

The selected examples make this limitation concrete without declaring the
reference wrong. In \texttt{mixed\_pln\_1278}, a pale field and the accepted
track look similar. In \texttt{mixed\_118}, dense woodland and leaf-covered
ground provide little visible support for the broad accepted foreground
region. In \texttt{mixed\_459}, several nearby vegetation--track contours
appear visually possible. The masks may encode knowledge unavailable
in the RGB frame, and visual ambiguity does not imply that both regions would
support a vehicle equally.

Because the dataset provides one released annotation per image and no
inter-annotator agreement, the thesis cannot estimate label variance, separate
model error from annotation uncertainty, or provide an upper bound on
RGB-only performance. A model that reproduces the labels perfectly would
reproduce their assumptions; it would not necessarily predict vehicle outcome
perfectly.

\section{Limitations}

\begin{enumerate}
  \item \textbf{One small, imbalanced dataset.} Training uses 1,002 images from
  three environments, with Power Line dominating. No environment-stratified
  result or geographically independent dataset is reported.

  \item \textbf{Sample-level split.} The deterministic hash prevents direct
  identifier overlap, but it does not explicitly group visually adjacent or
  near-duplicate frames. The degree of correlation between training and
  validation samples is not measured.

  \item \textbf{Annotation uncertainty is unquantified.} There are no repeated
  labels, annotator confidence scores, vehicle trials, or agreement statistics.

  \item \textbf{RGB cannot observe every physical constraint.} Support,
  deformability, water depth, slope, clearance, and wheel slip may require
  geometry, temporal evidence, or interaction.

  \item \textbf{Policy dependence.} Blue-only and blue+green are explicit
  vehicle policies. Results do not transfer numerically to blue+green+red,
  four-class CaT, or a different vehicle.

  \item \textbf{Repeated benchmark use.} Test images are excluded from
  optimisation and automatic checkpoint selection, but the project reports
  many test runs. The benchmark is not a single-use blinded evaluation.

  \item \textbf{Limited statistical scope.} The controlled ablation has one
  seed. The architecture table has three selected seeds on one fixed split and uses
  population standard deviation, not confidence intervals or hypothesis tests.

  \item \textbf{Architecture and pretraining are confounded.} PIDNet-S and
  DDRNet start randomly, while the leading systems use pretrained encoders.
  The table ranks complete recipes and cannot isolate causal architecture
  effects.

  \item \textbf{The convergence rule is shared, not model-specific
  optimisation.} The follow-up removes the obviously early 15-epoch stopping
  point for blue+green, but all models still share one objective, learning-rate
  schedule, and patience rule. Individually tuned recipes could change the
  ranking, and the blue-only matrix remains a fixed-budget result.

  \item \textbf{Two training software environments were used.} The FPN,
  U-Net, and SegFormer runs used PyTorch 2.13.0/CUDA 13.0; PIDNet, ROD,
  DDRNet, and BiSeNet used PyTorch 2.7.1/CUDA 12.6. The data, configuration,
  hardware partition, and evaluator were fixed, but framework version remains
  a cross-family confound.

  \item \textbf{Checkpoint availability differs from checkpoint identity.}
  The convergence suite records server-side SHA-256 identities for all 27
  selected checkpoints, but their binaries are not present locally. The
  reference PIDNet checkpoint remains the fully local reproduction chain.

  \item \textbf{Convergence training time is not usable evidence.} Shared-GPU
  workers were paused and resumed to control contention. This preserves
  training state but makes their recorded wall-clock duration unsuitable for
  comparing training cost.

  \item \textbf{Qualitative comparison is centred on PIDNet.} The detailed
  percentile, failure, and ALICE panels use the reference PIDNet model.
  Equivalent ROD and top-FPN failure panels were not produced, so qualitative
  error differences between architectures are not established.

  \item \textbf{ALICE is unlabelled.} The model output can be inspected, but no
  cross-domain accuracy or error rate can be calculated.

  \item \textbf{Pixel metrics omit planning and time.} No temporal consistency,
  confidence calibration, trajectory collision, intervention, or progress
  metric is evaluated.

  \item \textbf{Hardware is not the target embedded computer.} The Ryzen,
  H100 MIG partition, and desktop RTX~5060 do not measure automotive power,
  thermal limits, camera capture, planner latency, or end-to-end control.
\end{enumerate}

\section{Future Work}

Future work should close the largest evidence gaps before adding more
architectures.

\textbf{First, quantify the target.} A representative CaT subset should be
annotated independently by several people using a written vehicle-capability
guide. Agreement, boundary dispersion, and confidence should be reported.
Where possible, annotations should be paired with slope, roughness, soil, and
vehicle-trial evidence. This would separate visually ambiguous pixels from
clear model failures.

\textbf{Second, validate domain transfer.} A temporally coherent ALICE subset
should receive ground truth and vehicle-specific policy labels. Evaluation
should include sequence-level consistency and should hold complete locations
or runs out of training. ORFD, RUGD, RELLIS-3D, or WildScenes can extend
appearance diversity, but their labels must be mapped carefully rather than
treated as equivalent to CaT capability.

\textbf{Third, connect perception to planning.} The mask should be projected
into vehicle coordinates and evaluated with a planner. Metrics should weight
the near field, measure connected false-safe components, record whether the
planned footprint intersects blocked terrain, and quantify collision-free
progress. Obstacle inflation and a minimum-corridor rule can then be tested
against both exposure to blocked terrain and route availability.

\textbf{Fourth, represent uncertainty.} Probability calibration, ensembles,
test-time uncertainty, and an explicit unknown/abstain region should be
evaluated. Thresholds should be selected on validation data from a declared
application cost rather than by maximising test mIoU. Temporal smoothing should
be tested without allowing old confident errors to persist.

\textbf{Fifth, isolate modelling factors.} Pretrained and random initialisation
should be compared within the same encoder family. The selected architectures
should receive model-specific schedules and then be reevaluated under a common
compute budget. More seeds and grouped or location-held-out splits would give
stronger uncertainty estimates. Boundary supervision and hard-example sampling
should be assessed against FSR, FBR, and annotation disagreement.

\textbf{Finally, benchmark the actual robot stack.} The selected checkpoint
should run from camera capture through preprocessing, inference, projection,
planning, and control on the intended embedded computer. Measurements should
include median and tail latency, dropped frames, power, memory, temperature,
and throttling. RGB should be fused with depth, LiDAR, or vehicle interaction
where physical support cannot be inferred visually.

These steps would convert the present verified perception benchmark into
evidence about autonomous navigation rather than merely adding another
segmentation score.
